\chapter{Expanders and Ramanujan Graphs}

\begin{center}\itshape
A sparse graph in which every set has many neighbours, a random walk that forgets where it started in a few steps,\\
and a zeta function whose poles sit on a critical line:\\
three faces of one object, and the line is the one Riemann drew.
\end{center}

In 1973 two Soviet mathematicians, working on the same problem for the telephone network, took opposite roads. Mark Pinsker, an information theorist in Moscow, needed \emph{concentrators}: switching networks with few
wires in which any set of callers can be connected to the exchange by disjoint lines. He proved, at the 7th International Teletraffic Conference in Stockholm, that a random sparse graph does the job with
high probability. But a random graph is a promise, not a blueprint: nobody knew how to write one down. The same year Grigory Margulis, in Leningrad, published the first explicit family in the journal
\emph{Problemy Peredachi Informatsii}, built from the group $\mathrm{SL}_2(\Z)$ acting on the torus and from Kazhdan's property~(T), a notion from the representation theory of Lie groups that nobody
had expected to meet in a switchboard.

These graphs are now called \emph{expanders}, and Appendix~A of Volume~11 met them: a $d$-regular graph whose second eigenvalue is well below $d$ spreads its edges as if at random (the expander mixing lemma), and
the second eigenvalue cannot go below $2\sqrt{d-1}$ by much (Alon and Boppana). Graphs attaining this bound are called \emph{Ramanujan graphs}, because the first explicit families, built in 1988 by Alexander
Lubotzky, Ralph Phillips and Peter Sarnak and independently by Margulis, owe their eigenvalue bound to a conjecture of Ramanujan on the coefficients of modular forms, proved by Eichler for the case needed
and by Deligne in general.

The number $2\sqrt{d-1}$ is the one Kesten found in Chapter~8: the spectral radius of the $d$-regular tree, times $d$. This chapter explains why. A large regular graph without short cycles looks locally
like the tree, so its eigenvalues are distributed according to the spectral measure of the tree (Kesten and McKay); its closed non-backtracking walks are counted by a zeta function invented by Yasutaka Ihara
in 1966; and Toshikazu Sunada observed that the graph is Ramanujan exactly when this zeta function satisfies the Riemann hypothesis. For graphs, the Riemann hypothesis is a theorem, sometimes, and it
implies a prime number theorem with a square root error term.

How often does it hold? At a conference in Lausanne in the late 1980s, Noga Alon and Peter Sarnak made a bet: Sarnak thought that Ramanujan graphs are rare among all regular graphs, Alon that almost all
of them are Ramanujan. ``It was not extremely serious'', Alon recalled; ``we didn't even agree on what we are betting.'' In December 2024 Jiaoyang Huang, Theo McKenzie and Horng-Tzer Yau settled it with
random matrix theory: a large random regular graph is Ramanujan with probability about $69\%$ (Quanta Magazine, 18 April 2025). ``Both of us were somewhat wrong,'' said Alon, ``but I was a little bit more
correct, because the probability is bigger than half.''

\begin{reminder}[from Appendix A and Chapter 8]
For a $d$-regular graph $G$ on $n$ vertices with adjacency matrix $A$ and eigenvalues $d=\lambda_1\ge\lambda_2\ge\dots\ge\lambda_n$, $\lambda=\max(|\lambda_2|,|\lambda_n|)$; $G$ is then an \emph{$(n,d,\lambda)$-graph}.
Cheeger's inequality $h(G)\ge\frac{d-\lambda_2}2$ (Volume~11, Theorem~A.9); the expander mixing lemma $|e(S,T)-\frac dn|S||T||\le\lambda\sqrt{|S||T|}$ (Volume~11, Lemma~A.10); the Alon--Boppana bound
$\lambda_2\ge2\sqrt{d-1}-o(1)$ for graphs of large diameter (Volume~11, Theorem~A.11); the walk $p_{t+1}=\frac1dAp_t$ satisfies $\|p_t-u\|_2\le(\frac\lambda d)^t$, where $u$ is the uniform law (Volume~11, Exercise~A.8).
The $d$-regular tree $T_d$ has spectral radius $\frac{2\sqrt{d-1}}d$ (Theorem~\ref{thm:kesten} and Exercise~\ref{ex:regulartree}). A connected $d$-regular graph has $\lambda_1=d$ simple, and $\lambda_n=-d$ if and only if it is bipartite
(Exercise~\ref{ex:bipartiteeig}).
\end{reminder}

\section{Spectral gap and mixing}

Throughout, $G=(V,E)$ is a finite connected $d$-regular graph, $P=\frac1dA$ the transition matrix of simple random walk, and $\alpha=\frac\lambda d$. The matrix $P$ is symmetric, $P\mathbf 1=\mathbf 1$, and on the orthogonal of
$\mathbf 1$ its norm is $\alpha$. Write $J=\frac1n\mathbf 1\mathbf 1^{\mathsf T}$ for the orthogonal projection on the constants; then
\begin{equation}\label{eq:pje}
P=J+E,\qquad EJ=JE=0,\qquad\|E\|=\alpha.
\end{equation}

\begin{proposition}[mixing]\label{prop:mixing}
Whatever the initial law $p_0$, the law $p_t$ of the walk at time $t$ satisfies $\|p_t-u\|_{\mathrm{TV}}\le\frac{\sqrt n}2\alpha^t$. In particular, if $\alpha<1$ the walk is within $\varepsilon$ of uniform after
$t\ge\frac{\log(\sqrt n/(2\varepsilon))}{\log(1/\alpha)}$ steps.
\end{proposition}

\begin{proof}
By Volume~11, Exercise~A.8, $\|p_t-u\|_2\le\alpha^t\|p_0-u\|_2\le\alpha^t$, since $\|p_0-u\|_2^2=\|p_0\|_2^2-\frac1n\le1$. By Cauchy--Schwarz, $\|p_t-u\|_{\mathrm{TV}}=\frac12\|p_t-u\|_1\le\frac{\sqrt n}2\|p_t-u\|_2$.
\end{proof}

So on an expander family ($\alpha\le1-\varepsilon_0$ for all $n$) the walk mixes in $O(\log n)$ steps, the logarithm of the size: it is as fast as it could be, since after $t$ steps the walk has visited at most
$d^t$ vertices. On the cycle $C_n$ it needs about $n^2$ steps, and on the hypercube $\{0,1\}^k$ (with $n=2^k$, made lazy to avoid bipartiteness) about $k\log k=\log n\log\log n$.

\section{Saving random bits}

Expanders became a tool of computer science through the following observation, made by Mikl\'os Ajtai, J\'anos Koml\'os and Endre Szemer\'edi in 1987. Suppose a randomised algorithm uses $k$ random bits and errs with
probability at most $\frac12$, only on a set $B$ of ``bad'' seeds. Repeating it $t$ times with independent seeds reduces the error to $2^{-t}$, at the cost of $kt$ random bits. Instead, put the $n=2^k$ seeds on the vertices of a
$d$-regular expander, pick a uniform starting seed and walk: the seeds $X_0,X_1,\dots,X_t$ cost $k+t\log_2d$ bits. They are far from independent, yet they behave, for this purpose, almost as if they were.

\begin{proposition}[hitting property]\label{prop:hitting}
Let $B\subseteq V$ with $|B|=\beta n$, and let $X_0$ be uniform. Then
\[
P(X_0,X_1,\dots,X_t\in B)\le\beta(\beta+\alpha)^t.
\]
\end{proposition}

\begin{proof}
Let $\Pi$ be the diagonal matrix with $\Pi(x,x)=\mathbf 1_B(x)$. Summing over paths, $P(X_0,\dots,X_t\in B)=\frac1n\mathbf 1_B^{\mathsf T}(\Pi P\Pi)^t\mathbf 1_B$, because $\Pi^2=\Pi$. By \eqref{eq:pje}, $\Pi P\Pi=\Pi J\Pi+\Pi E\Pi$, where
$\Pi J\Pi=\frac1n\mathbf 1_B\mathbf 1_B^{\mathsf T}$ has norm $\frac{|B|}n=\beta$ and $\|\Pi E\Pi\|\le\alpha$. So $\|\Pi P\Pi\|\le\beta+\alpha$, and the probability is at most $\frac1n\|\mathbf 1_B\|^2(\beta+\alpha)^t=\beta(\beta+\alpha)^t$.
\end{proof}

With $\beta\le\frac12$ and $\alpha\le\frac14$ the error is at most $(\frac34)^t$, exponentially small, for $k+O(t)$ random bits instead of $kt$ (Exercise~\ref{ex:amplify}). For two-sided errors one needs more: the number of bad
seeds visited must be concentrated around its mean, as for independent samples. This is the Chernoff bound for expander walks, proved by David Gillman in 1993 (the journal version appeared in 1998).

\begin{theorem}[expander Chernoff bound]\label{thm:expchernoff}
Let $f:V\to[0,1]$ with mean $\mu=\frac1n\sum_xf(x)$, let $X_0$ be uniform and $S_t=\sum_{i=1}^tf(X_i)$. If $\alpha<1$, then for $0<\varepsilon\le1$,
\[
P\bigl(S_t\ge(\mu+\varepsilon)t\bigr)\le\exp\left(-\frac{(1-\alpha)\varepsilon^2t}{20}\right),
\]
and the same bound holds for $P(S_t\le(\mu-\varepsilon)t)$.
\end{theorem}

The bound is that of independent samples (Hoeffding's $e^{-2\varepsilon^2t}$, Volume~10) with a constant degraded by the spectral gap $1-\alpha$, and it does not depend on $n$.

\begin{proof}
Fix $r\in(0,1]$, let $D$ be the diagonal matrix with entries $e^{rf(x)}$, and let $w_k$ be the row vector $w_k=u^{\mathsf T}(PD)^k$, so that $w_k(x)=\mathbb E[e^{rS_k};X_k=x]$ and $\mathbb E[e^{rS_t}]=\sum_xw_t(x)$. Decompose each $w_k$ into
its mass $a_k=\sum_xw_k(x)$ and the rest: $w_k=a_ku^{\mathsf T}+w_k^\perp$ with $\sum_xw_k^\perp(x)=0$, and let $b_k=\sqrt n\,\|w_k^\perp\|_2$. Set $s=e^r-1$ and $m=\frac1n\sum_xe^{rf(x)}$.

\emph{One step.} Since $P$ is symmetric and doubly stochastic, $w_kP=a_ku^{\mathsf T}+z$, where $z=w_k^\perp P$ has mass $0$ and $\|z\|_2\le\alpha\|w_k^\perp\|_2$. Multiply by $D$.
\begin{itemize}
\item $a_ku^{\mathsf T}D$ has mass $a_km$, and its mass-free part has entries $\frac{a_k}n(e^{rf(x)}-m)$, all of absolute value at most $\frac{a_k}ns$ since $e^{rf(x)}$ and $m$ lie in $[1,e^r]$; so its norm is at most $\frac{a_ks}{\sqrt n}$.
\item $zD$ has mass $\sum_xz(x)(e^{rf(x)}-1)$, at most $\|z\|_2\sqrt n\,s$ in absolute value by Cauchy--Schwarz; its mass-free part has norm at most $\|zD\|_2\le e^r\|z\|_2$.
\end{itemize}
Hence, with $b_k=\sqrt n\|w_k^\perp\|_2$,
\[
a_{k+1}\le m\,a_k+s\alpha\,b_k,\qquad b_{k+1}\le s\,a_k+e^r\alpha\,b_k,\qquad a_0=1,\ b_0=0.
\]

\emph{Iterating.} Assume $r\le\frac{1-\alpha}4$. Then $s\le r+r^2\le\frac54r$ (as $e^r-1-r\le r^2\sum_{j\ge2}\frac1{j!}\le r^2$), so $e^r\alpha\le\alpha+\frac54r\alpha\le\frac{1+\alpha}2$, and since $m\ge1$, $m-e^r\alpha\ge\frac{1-\alpha}2$. Let
$\kappa=\frac{2s}{1-\alpha}$ and $\Lambda=m+s\alpha\kappa$. By induction, $a_k\le\Lambda^k$ and $b_k\le\kappa\Lambda^k$: the first from $m\Lambda^k+s\alpha\kappa\Lambda^k=\Lambda^{k+1}$, the second from
$s+e^r\alpha\kappa\le\kappa\Lambda$, which holds because $\kappa(\Lambda-e^r\alpha)\ge\kappa\frac{1-\alpha}2=s$.

\emph{Estimating $\Lambda$.} By convexity, $e^{rx}\le1+x(e^r-1)$ on $[0,1]$, so $m\le1+\mu s\le e^{\mu(r+r^2)}\le e^{\mu r+r^2}$. And $s\alpha\kappa=\frac{2s^2\alpha}{1-\alpha}\le\frac{4r^2}{1-\alpha}$ (as $s^2\le\frac{25}{16}r^2$). So
$\Lambda\le e^{\mu r+r^2}\bigl(1+\frac{4r^2}{1-\alpha}\bigr)\le\exp\bigl(\mu r+\frac{5r^2}{1-\alpha}\bigr)$.

\emph{Conclusion.} By Markov's inequality,
\[
P\bigl(S_t\ge(\mu+\varepsilon)t\bigr)\le e^{-r(\mu+\varepsilon)t}\,\mathbb E[e^{rS_t}]\le\exp\left(-r\varepsilon t+\frac{5r^2t}{1-\alpha}\right).
\]
Take $r=\frac{\varepsilon(1-\alpha)}{10}$, which is at most $\frac{1-\alpha}4$: the exponent is $-\frac{\varepsilon^2(1-\alpha)t}{10}+\frac{\varepsilon^2(1-\alpha)t}{20}=-\frac{\varepsilon^2(1-\alpha)t}{20}$. For the lower tail apply this to $1-f$, of mean $1-\mu$.
\end{proof}

\begin{think}
Why does the proof keep track of two numbers, the mass and the size of the rest, instead of bounding the norm of $PD$? (Hint: $\|D^{1/2}PD^{1/2}\|$ can be as large as $e^r$, which gives nothing; the point is that $P$
destroys the mass-free part at rate $\alpha$, while $D$ creates only a little of it, at rate $s\approx r$.)
\end{think}

\begin{remark}
The constant $20$ is not optimal. The sharpest form, due to Alexander Healy (2008) and Chung, Lam, Liu and Mitzenmacher (2012), gives $e^{-c(1-\alpha)\varepsilon^2t}$ with a better $c$ and extends to non-uniform starting laws
at the cost of a factor $\|p_0/u\|_2$. Expander walks are used to amplify randomised algorithms with few random bits (Cohen and Wigderson, Impagliazzo and Zuckerman, 1989), in pseudorandom generators for space-bounded
computation, and in Omer Reingold's 2005 proof that connectivity in undirected graphs can be decided in logarithmic space.
\end{remark}

\section{The spectral measure of the tree}

Why $2\sqrt{d-1}$? A large $d$-regular graph with few short cycles looks, around most vertices, like the $d$-regular tree $T_d$, so the number of closed walks of length $k$ at most vertices is the same as in the tree.
These numbers are the moments of the eigenvalue distribution. We first compute the tree's numbers, then identify the measure.

A walk $x_0,x_1,\dots,x_k$ is \emph{non-backtracking} if $x_{i+1}\ne x_{i-1}$ for $0<i<k$. Let $A_k(x,y)$ be the number of non-backtracking walks of length $k$ from $x$ to $y$.

\begin{lemma}[non-backtracking polynomials]\label{lem:nbpoly}
In a $d$-regular graph, finite or infinite, $A_0=I$, $A_1=A$, $A_2=A^2-dI$, and $A_{k+1}=AA_k-(d-1)A_{k-1}$ for $k\ge2$. Hence $A_k=q_k(A)$ for polynomials $q_k$ of degree $k$ with leading coefficient $1$, defined by $q_0=1$,
$q_1=x$, $q_2=x^2-d$ and $q_{k+1}=xq_k-(d-1)q_{k-1}$ for $k\ge2$.
\end{lemma}

\begin{proof}
$(A_kA)(x,y)$ counts the non-backtracking walks of length $k$ from $x$ to some $z$ followed by one step $z\to y$. Either this step does not backtrack, and the result is a non-backtracking walk of length $k+1$, or
$y$ is the vertex visited just before $z$. In the second case, removing the last two steps leaves a non-backtracking walk of length $k-1$ from $x$ to $y$, and conversely such a walk extends by a step $y\to z$ in $d-1$
non-backtracking ways if $k-1\ge1$, in $d$ ways if $k-1=0$. So $A_kA=A_{k+1}+(d-1)A_{k-1}$ for $k\ge2$ and $A_1A=A_2+dI$; and $A_kA=AA_k$ since $A_k$ is a polynomial in $A$ by induction.
\end{proof}

In a tree a non-backtracking walk never returns to its start (it is a geodesic), so at any vertex $o$ of $T_d$,
\begin{equation}\label{eq:treenb}
q_k(A)(o,o)=0\qquad(k\ge1).
\end{equation}

\begin{theorem}[Kesten, 1959]\label{thm:kestenmeasure}
Let $d\ge3$ and let $a_k$ be the number of closed walks of length $k$ from a vertex $o$ of the $d$-regular tree. Then $a_k=\int x^k\nu_d(\dd x)$, where $\nu_d$ is the probability measure
\[
\nu_d(\dd x)=\frac{d\sqrt{4(d-1)-x^2}}{2\pi(d^2-x^2)}\,\mathbf 1_{|x|\le2\sqrt{d-1}}\dd x.
\]
\end{theorem}

\begin{proof}
\emph{Uniqueness.} Let $L$ be the linear form on polynomials with $L(x^k)=a_k=A^k(o,o)$. By \eqref{eq:treenb}, $L(q_k)=0$ for $k\ge1$ and $L(q_0)=1$. Since $q_k=x^k+$ (lower terms), these equations determine $L(x^k)$ by
induction on $k$. So it suffices to show that $\int q_k\dd\nu_d=0$ for $k\ge1$ and $\int\dd\nu_d=1$.

\emph{Change of variables.} Let $\sigma=\sqrt{d-1}$ and $x=2\sigma\cos\theta$, $\theta\in[0,\pi]$. Then $\sqrt{4(d-1)-x^2}=2\sigma\sin\theta$ and $d^2-x^2=d^2-4\sigma^2\cos^2\theta=\sigma^4|1-\sigma^{-2}e^{2i\theta}|^2$ (expand:
$\sigma^4(1+\sigma^{-4}-2\sigma^{-2}\cos2\theta)=\sigma^4+1-2\sigma^2(2\cos^2\theta-1)=(\sigma^2+1)^2-4\sigma^2\cos^2\theta$, and $\sigma^2+1=d$). With $\rho=\sigma^{-2}=\frac1{d-1}<1$,
\[
\nu_d(\dd x)=\frac{2d}{\pi(d-1)}\,\frac{\sin^2\theta}{|1-\rho e^{2i\theta}|^2}\dd\theta.
\]

\emph{Chebyshev form of $q_k$.} Let $U_k$ be the Chebyshev polynomials of the second kind, $U_k(\cos\theta)=\frac{\sin(k+1)\theta}{\sin\theta}$, with $U_{-1}=0$; they satisfy $U_{k+1}=2xU_k-U_{k-1}$. One checks by induction that
$q_k(2\sigma\cos\theta)=\sigma^k\bigl(U_k-\rho U_{k-2}\bigr)(\cos\theta)$ for $k\ge1$: for $k=1,2$ this is $2\sigma\cos\theta$ and $\sigma^2(4\cos^2\theta-1-\rho)=4\sigma^2\cos^2\theta-d$; and both sides satisfy the recursion
$q_{k+1}=xq_k-\sigma^2q_{k-1}$ for $k\ge2$, the right side because $\sigma^{k+1}(2\cos\theta\,V_k-V_{k-1})=\sigma^{k+1}V_{k+1}$ for $V_k=U_k-\rho U_{k-2}$. Hence
\[
\sin\theta\,q_k(x)=\sigma^k\bigl(\sin(k+1)\theta-\rho\sin(k-1)\theta\bigr)=\sigma^k\operatorname{Im}\bigl(e^{i(k-1)\theta}(e^{2i\theta}-\rho)\bigr).
\]

\emph{The integrals.} Since $|1-\rho e^{2i\theta}|=|e^{2i\theta}-\rho|$, we have $\frac{e^{2i\theta}-\rho}{|1-\rho e^{2i\theta}|^2}=\frac1{e^{-2i\theta}-\rho}=\sum_{j\ge0}\rho^je^{2i(j+1)\theta}$, uniformly convergent. So for $k\ge1$
\[
\int q_k\dd\nu_d=\frac{2d\sigma^k}{\pi(d-1)}\sum_{j\ge0}\rho^j\int_0^\pi\sin\theta\,\sin\bigl((k+1+2j)\theta\bigr)\dd\theta=0 ,
\]
because $\int_0^\pi\sin\theta\sin m\theta\dd\theta=0$ for every integer $m\ge2$. For the mass, $\frac1{|1-\rho e^{2i\theta}|^2}=\frac1{1-\rho^2}\sum_{j\in\Z}\rho^{|j|}e^{2ij\theta}$ (the Poisson kernel), and $\sin^2\theta=\frac12-\frac12\cos2\theta$, so
\[
\int_0^\pi\frac{\sin^2\theta\dd\theta}{|1-\rho e^{2i\theta}|^2}=\frac1{1-\rho^2}\Bigl(\frac\pi2-\frac\pi4\cdot2\rho\Bigr)=\frac\pi{2(1+\rho)} ,
\]
and $\nu_d$ has mass $\frac{2d}{\pi(d-1)}\cdot\frac{\pi(d-1)}{2d}=1$, since $1+\rho=\frac d{d-1}$.
\end{proof}

The measure $\nu_d$ is the \emph{Kesten--McKay law}. Its support is $[-2\sqrt{d-1},2\sqrt{d-1}]$, so the return probabilities $p_{2n}=a_{2n}/d^{2n}$ of simple random walk satisfy $p_{2n}^{1/(2n)}\to\frac{2\sqrt{d-1}}d$: Kesten's
theorem again, now as the edge of a spectrum. As $d\to\infty$, rescaled by $\sqrt{d-1}$, it becomes Wigner's semicircle law of the next chapter (Exercise~\ref{ex:kmsemicircle}).

\begin{theorem}[McKay, 1981]\label{thm:mckay}
Let $G_N$ be $d$-regular graphs on $N\to\infty$ vertices, $d\ge3$, such that for every $r$ the proportion of vertices $x$ whose ball $B(x,r)$ contains a cycle tends to $0$. Then the empirical distribution of the
eigenvalues of $A(G_N)$, $\frac1N\sum_i\delta_{\lambda_i}$, converges weakly to $\nu_d$.
\end{theorem}

\begin{proof}
Fix $k$ and $r\ge k$. If $B(x,r)$ contains no cycle, the subgraph it induces is a tree in which every vertex at distance less than $r$ from $x$ has degree $d$, so it is isomorphic to the ball of radius $r$ in $T_d$. A
closed walk of length $k$ from $x$ stays at distance at most $\frac k2$ from $x$ and uses only edges between vertices of $B(x,r)$; so the number of such walks is $a_k$. At the other vertices it is at most $d^k$. Hence
\[
\Bigl|\frac1N\sum_i\lambda_i^k-a_k\Bigr|=\Bigl|\frac1N\operatorname{tr}A^k-a_k\Bigr|\le\frac{\#\{x:B(x,r)\text{ has a cycle}\}}N\,2d^k\longrightarrow0.
\]
So all moments of the empirical distributions converge to those of $\nu_d$ (Theorem~\ref{thm:kestenmeasure}). All these measures live on $[-d,d]$, where every continuous function is a uniform limit of polynomials
(Weierstrass), so $\int g\dd\mu_N\to\int g\dd\nu_d$ for every continuous $g$.
\end{proof}

The hypothesis holds for a random $d$-regular graph with high probability: its number of cycles of length $\ell$ converges in law to a Poisson variable of mean $\frac{(d-1)^\ell}{2\ell}$ (Bollob\'as, 1980), so only $O(1)$ vertices
are near a short cycle. Figure~\ref{fig:kestenmckay} shows the agreement for $N=3000$.

\begin{corollary}\label{cor:serre}
Under the hypothesis of Theorem~\ref{thm:mckay}, for every $\varepsilon>0$ a proportion at least $c(\varepsilon,d)>0$ of the eigenvalues of $G_N$ are at least $2\sqrt{d-1}-\varepsilon$, for $N$ large.
\end{corollary}

\begin{proof}
Let $g$ be continuous, $0\le g\le1$, zero on $(-\infty,2\sqrt{d-1}-\varepsilon]$ and equal to $1$ on $[2\sqrt{d-1}-\frac\varepsilon2,\infty)$. Then $\int g\dd\nu_d>0$ since the density of $\nu_d$ is positive inside its support, and the
proportion of eigenvalues above $2\sqrt{d-1}-\varepsilon$ is at least $\frac1N\sum_ig(\lambda_i)\to\int g\dd\nu_d$.
\end{proof}

So $2\sqrt{d-1}$ is not only a lower bound for $\lambda_2$ (Alon--Boppana): a positive fraction of the whole spectrum crowds up to it. Jean-Pierre Serre proved in 1997 a version for all regular graphs, without the
tree-like hypothesis. Ramanujan graphs are the graphs in which nothing at all sticks out of the tree's spectrum.

\begin{figure}[ht]
\centering
\fig{ch09-kestenmckay}
\caption{Left: histogram of the $3000$ eigenvalues of a random $4$-regular graph and the Kesten--McKay density $\nu_4$ (red), supported on $[-2\sqrt3,2\sqrt3]$; the trivial eigenvalue $4$ is the isolated dot. Right: the
poles $u$ of the Ihara zeta function of a random $3$-regular graph on $300$ vertices (Section~\ref{sec:ihara}). The trivial poles $u=1$ and $u=\frac12$ (orange) lie on the real axis, and so do the poles $\pm1$ of the factor $(1-u^2)^{m-n}$ (not drawn); all the others lie on the circle
$|u|=\frac1{\sqrt2}$: this graph is Ramanujan, its zeta function satisfies the Riemann hypothesis.}
\label{fig:kestenmckay}
\end{figure}

\section{Ihara's zeta function}\label{sec:ihara}

In 1966 Yasutaka Ihara, studying discrete subgroups of $\mathrm{SL}_2$ over $p$-adic fields, defined a zeta function that counts closed geodesics, as Selberg's zeta function does on a hyperbolic surface. Jean-Pierre Serre
remarked in his book \emph{Trees} (1977) that it makes sense for any finite graph, Toshikazu Sunada developed this in 1986, and Ki-ichiro Hashimoto (1989) and Hyman Bass (1992) proved the determinant formulas
that we now prove for regular graphs.

Let $G$ be a finite connected simple graph with $n$ vertices and $m$ edges. Each edge gives two \emph{oriented edges}; for an oriented edge $e$ from $x$ to $y$ write $o(e)=x$, $t(e)=y$, and $\bar e$ for the reversed edge.
The \emph{non-backtracking} (or \emph{Hashimoto}) matrix is the $2m\times2m$ matrix
\[
B(e,f)=1\ \text{ if }t(e)=o(f)\text{ and }f\ne\bar e,\qquad B(e,f)=0\ \text{ otherwise.}
\]
A \emph{cycle} of length $k$ is a sequence $c=(e_1,\dots,e_k)$ of oriented edges with $B(e_i,e_{i+1})=1$ for all $i$, indices modulo $k$: a closed walk that never backtracks, not even from its last edge to its first.
Their number is $N_k=\operatorname{tr}B^k$. Two cycles are \emph{equivalent} if one is a cyclic rotation of the other, and a cycle is \emph{primitive} if it is not a power $P^j$, $j\ge2$, of a shorter cycle $P$ (the
cycle $P$ run $j$ times). Primitive cycles play the role of primes: let $\pi(k)$ be the number of equivalence classes of primitive cycles of length $k$, and define the \emph{Ihara zeta function}
\[
Z_G(u)=\prod_{[P]}\bigl(1-u^{\ell(P)}\bigr)^{-1}=\prod_{k\ge1}(1-u^k)^{-\pi(k)},
\]
the product over the classes of primitive cycles, $\ell(P)$ being the length.

\begin{proposition}\label{prop:iharadet}
For $|u|$ small enough, $\displaystyle\log Z_G(u)=\sum_{k\ge1}\frac{N_k}ku^k$ and $Z_G(u)=\det(I-uB)^{-1}$. In particular $Z_G$ is the inverse of a polynomial.
\end{proposition}

\begin{proof}
If $c$ is a cycle of length $k$, the rotations that fix it form a subgroup of $\Z/k$, generated by the rotation by some $\ell\mid k$; then $c=P^{k/\ell}$, where $P$ is made of the first $\ell$ edges of $c$, and $P$ is primitive.
Conversely, a primitive class of length $\ell$ consists of exactly $\ell$ distinct cycles (no rotation fixes a primitive cycle), and each gives one cycle of length $k$ for each multiple $k$ of $\ell$. Hence
\begin{equation}\label{eq:nkpi}
N_k=\sum_{\ell\mid k}\ell\,\pi(\ell).
\end{equation}
Then
\[
\log Z_G(u)=\sum_{[P]}\sum_{j\ge1}\frac{u^{j\ell(P)}}j=\sum_k\frac{u^k}k\sum_{\ell\mid k}\ell\pi(\ell)=\sum_k\frac{N_k}ku^k ,
\]
all series converging absolutely for $|u|<1/b$, where $b$ is the maximal row sum of $B$, since
$N_k\le2mb^{k}$. Finally, if $\mu_1,\dots,\mu_{2m}$ are the eigenvalues of $B$ (with $|\mu_i|\le b$), then $\sum_k\frac1k\operatorname{tr}B^ku^k=\sum_i\sum_k\frac1k(\mu_iu)^k$, which equals $-\sum_i\log(1-\mu_iu)=-\log\det(I-uB)$.
\end{proof}

\begin{theorem}[Ihara, 1966; Bass, 1992]\label{thm:bass}
If $G$ is $(q+1)$-regular,
\[
\det(I-uB)=(1-u^2)^{m-n}\det\bigl(I-uA+qu^2I\bigr).
\]
\end{theorem}

\begin{proof}
Let $S$ and $T$ be the $n\times2m$ matrices $S(x,e)=\mathbf 1_{o(e)=x}$ and $T(x,e)=\mathbf 1_{t(e)=x}$, and $R$ the $2m\times2m$ permutation matrix $R(e,f)=\mathbf 1_{f=\bar e}$. Directly from the definitions,
\[
ST^{\mathsf T}=A,\qquad TT^{\mathsf T}=(q+1)I,\qquad SR=T,\qquad B=T^{\mathsf T}S-R ,
\]
the last because $(T^{\mathsf T}S)(e,f)=\mathbf 1_{t(e)=o(f)}$ and $\bar e$ is one of the $f$ with $o(f)=t(e)$. Let $M=I+uR$. Since $R^2=I$ and $R$ swaps the two orientations of each edge, $M$ is made of $m$ blocks
$\bigl(\begin{smallmatrix}1&u\\u&1\end{smallmatrix}\bigr)$, so $\det M=(1-u^2)^m$, and $M^{-1}=\frac{I-uR}{1-u^2}$ for $u^2\ne1$. By Sylvester's identity $\det(I-XY)=\det(I-YX)$ (proved by computing the determinant of
$\bigl(\begin{smallmatrix}I&X\\Y&I\end{smallmatrix}\bigr)$ by eliminating either block),
with $W=SM^{-1}T^{\mathsf T}$,
\begin{align*}
\det(I-uB)&=\det\bigl(M-uT^{\mathsf T}S\bigr)\\
&=\det M\cdot\det\bigl(I_{2m}-uM^{-1}T^{\mathsf T}S\bigr)\\
&=\det M\cdot\det(I_n-uW).
\end{align*}
Now $W=\frac{ST^{\mathsf T}-uSRT^{\mathsf T}}{1-u^2}=\frac{A-u(q+1)I}{1-u^2}$, so $(1-u^2)^n\det(I_n-uW)$ equals
\[
\det\bigl((1-u^2)I-uA+(q+1)u^2I\bigr)=\det(I-uA+qu^2I).
\]
This proves the identity for $u^2\ne1$; both sides are polynomials, so it holds for all $u$.
\end{proof}

So the zeta function of a regular graph is determined by the spectrum of $A$. Each eigenvalue $\lambda_j$ of $A$ contributes the factor $1-\lambda_ju+qu^2=(1-\mu_ju)(1-\mu'_ju)$, where $\mu_j,\mu'_j$ are the roots of $\mu^2-\lambda_j\mu+q$, and
comparing the two sides of Theorem~\ref{thm:bass} as polynomials in $u$ shows that the $2m$ eigenvalues of $B$ are the $\mu_j,\mu'_j$ together with $1$ and $-1$, each $m-n$ times. Therefore
\begin{equation}\label{eq:nkspec}
N_k=\sum_{j=1}^n\bigl(\mu_j^k+\mu_j'^k\bigr)+(m-n)\bigl(1+(-1)^k\bigr).
\end{equation}

\section{The Riemann hypothesis for graphs}

A connected $(q+1)$-regular graph is \emph{Ramanujan} if every eigenvalue $\lambda$ of $A$ with $|\lambda|\ne q+1$ satisfies $|\lambda|\le2\sqrt q$. In the variable $u=q^{-s}$, the trivial eigenvalue $\lambda_1=q+1$ gives
$1-(q+1)u+qu^2=(1-u)(1-qu)$, that is poles of $Z_G$ at $u=1$ and $u=\frac1q$, or $s=0$ and $s=1$: they play the role of the pole of $\zeta$ at $s=1$. The critical line $\operatorname{Re}s=\frac12$ is the circle $|u|=q^{-1/2}$.

\begin{theorem}[Sunada, 1986]\label{thm:graphRH}
Let $G$ be connected and $(q+1)$-regular, $q\ge2$. Then $G$ is Ramanujan if and only if every pole $u$ of $Z_G$ with $\frac1q<|u|<1$ satisfies $|u|=\frac1{\sqrt q}$; in the variable $s$, every pole with $0<\operatorname{Re}s<1$ lies on
the line $\operatorname{Re}s=\frac12$.
\end{theorem}

\begin{proof}
By Theorem~\ref{thm:bass} the poles of $Z_G$ are $\pm1$ and the roots of $qu^2-\lambda u+1$ for the eigenvalues $\lambda$ of $A$; these roots have product $\frac1q$. For $\lambda=\pm(q+1)$ they are $\pm1$ and $\pm\frac1q$, outside the open annulus.
If $|\lambda|\le2\sqrt q$, the discriminant $\lambda^2-4q$ is $\le0$, so the roots are complex conjugate (or equal), with $|u|^2=\frac1q$. If $2\sqrt q<|\lambda|<q+1$, the roots are real, distinct, of the sign of $\lambda$, with
product $\frac1q$, so one has modulus less than $\frac1{\sqrt q}$ and one more; both lie strictly between $\frac1q$ and $1$ in modulus, because the polynomial $qu^2-|\lambda|u+1$ is positive at $u=\frac1q$ and $u=1$
(its values there are $\frac{q+1-|\lambda|}q$ and $q+1-|\lambda|$) while its minimum, at $u=\frac{|\lambda|}{2q}\in(\frac1q,1)$, is negative. So a non-trivial eigenvalue outside $[-2\sqrt q,2\sqrt q]$ produces a pole off the circle, and
conversely.
\end{proof}

For the Riemann zeta function, the hypothesis implies the prime number theorem with error $O(\sqrt x\log x)$ (Volume~5). For graphs, the same causes have the same effect, and the proof is two lines.

\begin{theorem}[prime number theorem for graphs]\label{thm:graphpnt}
Let $G$ be a connected non-bipartite Ramanujan $(q+1)$-regular graph on $n$ vertices, $q\ge2$. For every $k\ge2$,
\[
\Bigl|\pi(k)-\frac{q^k}k\Bigr|\le\frac{6n\,q^{k/2}}k.
\]
\end{theorem}

\begin{proof}
In \eqref{eq:nkspec}, $\lambda_1=q+1$ gives $\mu=q$ and $\mu'=1$. For $j\ge2$, $|\lambda_j|\le2\sqrt q$ ($-(q+1)$ is not an eigenvalue since $G$ is not bipartite), so $|\mu_j|=|\mu'_j|=\sqrt q$ as in the previous proof. And
$m-n=\frac{n(q-1)}2$. So $|N_k-q^k|\le1+2(n-1)q^{k/2}+n(q-1)\le3nq^{k/2}$, using $q\le q^{k/2}$ for $k\ge2$. By \eqref{eq:nkpi}, $k\pi(k)=N_k-\sum_{\ell\mid k,\,\ell<k}\ell\pi(\ell)$, and the subtracted sum is between $0$ and
$\sum_{\ell\le k/2}N_\ell$. A cycle of length $\ell$ is a choice of a first oriented edge and then of $\ell-1$ non-backtracking steps, so $N_\ell\le2m\,q^{\ell-1}$, and
$\sum_{\ell\le k/2}N_\ell\le2m\frac{q^{k/2}}{q-1}=\frac{n(q+1)}{q-1}q^{k/2}\le3nq^{k/2}$. Together, $|k\pi(k)-q^k|\le6nq^{k/2}$.
\end{proof}

\begin{beauty}
In Volume~5 the explicit formula expressed the primes through the zeros of $\zeta$; here \eqref{eq:nkspec} expresses the closed geodesics of a graph through the eigenvalues of its adjacency matrix, and the Riemann hypothesis
becomes a statement about a finite symmetric matrix: no eigenvalue may stick out of the spectrum of the universal cover, the tree. Selberg's trace formula does the same for hyperbolic surfaces, with the Laplacian
in place of $A$ and the hyperbolic plane in place of the tree; the graph is the version where everything can be proved and computed.
\end{beauty}

\section{Ramanujan graphs}

By Alon--Boppana (Volume~11, Theorem~A.11) and Corollary~\ref{cor:serre}, a family of $(q+1)$-regular graphs with $n\to\infty$ cannot have $\lambda\le2\sqrt q-\varepsilon$: Ramanujan graphs are the best possible expanders. Do they exist,
for every degree and infinitely many sizes? The answer is a tour of twentieth century mathematics, and we only describe it.

\begin{remark}[explicit constructions]
For primes $p\ne\ell$, both congruent to $1$ modulo $4$, Lubotzky, Phillips and Sarnak (1988) consider the $p+1$ solutions of $a_0^2+a_1^2+a_2^2+a_3^2=p$ with $a_0>0$ odd and $a_1,a_2,a_3$ even (Jacobi's four square theorem counts them),
map each to a matrix of $\mathrm{PGL}_2(\mathbf F_\ell)$, and take the Cayley graph. The resulting $(p+1)$-regular graphs are Ramanujan: their eigenvalues are expressed through the Hecke eigenvalues of modular forms of
weight~$2$ attached to a quaternion algebra, and the bound $2\sqrt p$ is Ramanujan's conjecture for them, proved by Eichler (1954) using Weil's Riemann hypothesis for curves over finite fields. They also have girth at
least $\frac43\log_p\ell$, larger than any previous construction. Margulis found the same graphs independently (1988), and Moshe Morgenstern (1994) extended the construction to every degree $q+1$ with $q$ a prime power.
\end{remark}

\begin{remark}[existence for all degrees, and random graphs]
In 2013 Adam Marcus, Daniel Spielman and Nikhil Srivastava proved that bipartite Ramanujan graphs exist for every degree $d\ge3$ and infinitely many sizes, by showing that some signing of the adjacency matrix of a
Ramanujan graph has small eigenvalues, with their method of \emph{interlacing families} of polynomials; the same method solved the Kadison--Singer problem a few months later. Joel Friedman proved in 2008 that a random
$d$-regular graph has $\lambda\le2\sqrt{d-1}+\varepsilon$ with probability tending to $1$, as Alon had conjectured, and Huang, McKenzie and Yau proved in 2024 that the fluctuations of $\lambda_2$ around $2\sqrt{d-1}$ follow
the Tracy--Widom law of random matrix theory, at the scale $N^{-2/3}$, whence the probability of about $69\%$ of the opening story.
\end{remark}

\begin{remark}[Cayley graphs of $\mathrm{SL}_2(\mathbf F_p)$]
Let $S=\{(\begin{smallmatrix}1&\pm1\\0&1\end{smallmatrix}),(\begin{smallmatrix}1&0\\\pm1&1\end{smallmatrix})\}$. The Cayley graphs of $\mathrm{SL}_2(\mathbf F_p)$ with respect to $S$ (reduced modulo $p$) are $4$-regular, with $p(p^2-1)$
vertices (Exercise~\ref{ex:sl2gen}), and they form an expander family: this follows from Selberg's theorem of 1965 that the Laplacian on the modular curves $\Gamma(p)\backslash\mathbb H$ has no eigenvalue in $(0,\frac3{16})$,
by a transfer from surfaces to graphs explained in Lubotzky's book (1994). Jean Bourgain and Alex Gamburd proved in 2008 that the same holds for any set $S$ of integer matrices generating a Zariski dense subgroup of
$\mathrm{SL}_2$, with no automorphic forms at all, by additive combinatorics. These graphs are not Ramanujan in general, but their spectral gap is uniform in $p$; a direct computation gives it.
\end{remark}

\section{Exercises}

\begin{exercise}[$\star$]\label{ex:bipartiteeig}
Let $G$ be connected and $d$-regular. Show that $-d$ is an eigenvalue of $A$ if and only if $G$ is bipartite, and that then simple random walk does not converge to the uniform law. Show that the \emph{lazy} walk
$\frac12(I+P)$ has all its non-trivial eigenvalues in $[0,\frac{1+\lambda_2/d}2]$.
\end{exercise}

\begin{exercise}[$\star\star$]\label{ex:amplify}
A randomised algorithm reads $k$ random bits; on a ``yes'' instance it answers ``yes'' for at least half of the $2^k$ seeds, on a ``no'' instance it always answers ``no''. Using a $d$-regular graph on the $2^k$ seeds with
$\alpha\le\frac14$, design a procedure that uses $k+t\log_2d$ random bits and errs with probability at most $(\frac34)^t$.
\end{exercise}

\begin{exercise}[$\star$]\label{ex:treemoments}
Show that the number of closed walks of length $4$ from a vertex of $T_d$ is $d(2d-1)$, first by counting, then from $\int q_4\dd\nu_d=0$. Show also that $a_6=d(5d^2-6d+2)$.
\end{exercise}

\begin{exercise}[$\star\star$]\label{ex:kmsemicircle}
Show that the image of $\nu_d$ under $x\mapsto x/\sqrt{d-1}$ converges, as $d\to\infty$, to the semicircle law $\frac1{2\pi}\sqrt{4-y^2}\,\mathbf 1_{|y|\le2}\dd y$, and that $a_{2k}/(d-1)^k\to C_k$, the Catalan number.
\end{exercise}

\begin{exercise}[$\star$]\label{ex:cycleihara}
Show that the Ihara zeta function of the cycle $C_n$ is $(1-u^n)^{-2}$, directly and from Theorem~\ref{thm:bass} with $q=1$.
\end{exercise}

\begin{exercise}[$\star\star$]\label{ex:k4}
Show that the complete graph $K_4$ is Ramanujan. Compute the eigenvalues of its matrix $B$ and deduce $N_3=24$ and $N_4=24$; check that $\pi(3)=8$ and $\pi(4)=6$ by drawing the cycles.
\end{exercise}

\begin{exercise}[$\star\star$]\label{ex:irregularbass}
Let $G$ be any finite connected simple graph, and $Q$ the diagonal matrix with $Q(x,x)=\deg x-1$. Show that $\det(I-uB)=(1-u^2)^{m-n}\det(I-uA+u^2Q)$.
\end{exercise}

\begin{exercise}[$\star\star$]\label{ex:sl2gen}
Let $p$ be a prime. Show that $|\mathrm{SL}_2(\mathbf F_p)|=p(p^2-1)$ and that the matrices $U=(\begin{smallmatrix}1&1\\0&1\end{smallmatrix})$ and $L=(\begin{smallmatrix}1&0\\1&1\end{smallmatrix})$ generate $\mathrm{SL}_2(\mathbf F_p)$. Deduce that the Cayley
graph of the remark above is connected and $4$-regular for $p\ge3$.
\end{exercise}

\section{Hints}

\begin{hint}{ex:bipartiteeig}
If $Af=-df$, look at a vertex where $|f|$ is maximal. For a bipartite graph, $f=\pm1$ on the two sides.
\end{hint}

\begin{hint}{ex:amplify}
Run the algorithm on the seeds $X_0,\dots,X_t$ of a walk and answer ``yes'' if one run does. Apply Proposition~\ref{prop:hitting} to the set of seeds answering ``no''.
\end{hint}

\begin{hint}{ex:treemoments}
A closed walk of length $4$ either goes out and back twice, or goes out two steps and comes back. Compute $q_3$ and $q_4$ from the recursion.
\end{hint}

\begin{hint}{ex:kmsemicircle}
Write the density of the image measure explicitly. For the moments, use dominated convergence on the compact support.
\end{hint}

\begin{hint}{ex:cycleihara}
There are exactly two primitive classes. For Bass's formula, the eigenvalues of $C_n$ are $2\cos\frac{2\pi j}n$ and $1-2u\cos\theta+u^2=|1-ue^{i\theta}|^2$.
\end{hint}

\begin{hint}{ex:k4}
The eigenvalues of $K_4$ are $3$ and $-1$ (three times). For $\lambda=-1$, $\mu^2+\mu+2=0$, and $\mu^3=2-\mu$, $\mu^4=3\mu+2$.
\end{hint}

\begin{hint}{ex:irregularbass}
Repeat the proof of Theorem~\ref{thm:bass}: only $TT^{\mathsf T}$ changes.
\end{hint}

\begin{hint}{ex:sl2gen}
Count the first column, then the second. For generation, $U^a$ and $L^b$ perform the row operations of Gaussian elimination.
\end{hint}

\section{Solutions}

\begin{solution}{ex:bipartiteeig}
\idea An eigenvector for $-d$ must flip sign along every edge, which is exactly a two-colouring.
\writeup If $Af=-df$ with $f\ne0$, let $x$ maximise $|f|$, say $f(x)=M>0$. Then $\sum_{y\sim x}f(y)=-dM$ with each $f(y)\ge-M$, so $f(y)=-M$ for all neighbours, and then $f=M$ on their neighbours, and so on: by connectedness
$f=\pm M$ everywhere and adjacent vertices have opposite signs, so $\{f>0\}$, $\{f<0\}$ is a bipartition. Conversely, for a bipartition $V_1\cup V_2$, $f=\mathbf 1_{V_1}-\mathbf 1_{V_2}$ satisfies $Af=-df$. Then the walk alternates sides, so
$p_t$ is supported on one side at each time and cannot converge to $u$. The lazy walk has eigenvalues $\frac{1+\lambda_i/d}2\in[0,1]$, and the largest non-trivial one is $\frac{1+\lambda_2/d}2<1$.
\end{solution}

\begin{solution}{ex:amplify}
\idea A walk that stays in the bad set for $t+1$ steps is as unlikely as $t$ independent failures, up to the spectral gap.
\writeup Choose $X_0$ uniform ($k$ bits) and $t$ steps of the walk ($\log_2d$ bits each, with $d$ a power of $2$, or $\lceil\log_2d\rceil$ bits in general), run the algorithm on each seed, and answer ``yes'' if some run does. On a ``no''
instance the answer is always right. On a ``yes'' instance the bad set $B$ of seeds answering ``no'' has $\beta\le\frac12$, and by Proposition~\ref{prop:hitting} the error probability is
$P(X_0,\dots,X_t\in B)\le\beta(\beta+\alpha)^t\le\frac12(\frac34)^t$.
\end{solution}

\begin{solution}{ex:treemoments}
\idea Count by shapes; the polynomial identity gives the same number with no counting at all.
\writeup Out and back twice: $d\cdot d$ ways. Out two steps along a path and back: $d(d-1)$ ways. Total $d^2+d(d-1)=d(2d-1)$. From the recursion, $q_3=x^3-(2d-1)x$ and $q_4=xq_3-(d-1)q_2=x^4-(3d-2)x^2+d(d-1)$; since
$a_2=d$, $\int q_4\dd\nu_d=0$ gives $a_4=(3d-2)d-d(d-1)=2d^2-d$. Similarly $q_5=xq_4-(d-1)q_3=x^5-(4d-3)x^3+(3d^2-4d+1)x$ and $q_6=xq_5-(d-1)q_4=x^6-(5d-4)x^4+(6d^2-9d+3)x^2-d(d-1)^2$, so
$a_6=(5d-4)a_4-(6d^2-9d+3)a_2+d(d-1)^2=(5d-4)(2d^2-d)-d(6d^2-9d+3)+d(d-1)^2=d(5d^2-6d+2)$. For $d=3$: $a_4=15$, $a_6=87$.
\end{solution}

\begin{solution}{ex:kmsemicircle}
\idea For large $d$ the factor $\frac{d^2}{d^2-x^2}$ is close to $1$ on the support, and what remains is a semicircle.
\writeup With $x=\sqrt{d-1}\,y$, the image has density $\frac{d(d-1)\sqrt{4-y^2}}{2\pi(d^2-(d-1)y^2)}\mathbf 1_{|y|\le2}$, and $\frac{d(d-1)}{d^2-(d-1)y^2}\to1$ uniformly for $|y|\le2$ (the denominator is at least $d^2-4(d-1)=(d-2)^2$). So the
densities converge uniformly on $[-2,2]$, hence the measures converge, and so do the moments: $\frac{a_{2k}}{(d-1)^k}=\int y^{2k}\to\frac1{2\pi}\int_{-2}^2y^{2k}\sqrt{4-y^2}\dd y=C_k$ (the moments of the semicircle are the Catalan numbers, Chapter~10). Combinatorially:
for large $d$, the closed walks of length $2k$ that dominate are those which never step back except along the edge they came from, and they are counted by Dyck paths, $C_k$ shapes times about $d^k$ choices.
\end{solution}

\begin{solution}{ex:cycleihara}
\idea The only geodesics of a cycle go round it, in one of two directions.
\writeup A non-backtracking walk on $C_n$ keeps going in the same direction, so a cycle goes round $j$ times in one direction, with length $jn$; the primitive classes are the two loops of length $n$. So $Z=(1-u^n)^{-2}$. With
Theorem~\ref{thm:bass}, $q=1$ and $m=n$: $\det(I-uB)=\prod_j(1-2u\cos\frac{2\pi j}n+u^2)=\prod_j(1-ue^{2\pi ij/n})(1-ue^{-2\pi ij/n})=(1-u^n)^2$, since $\prod_j(1-u\omega^j)=1-u^n$ over the $n$-th roots of unity $\omega^j$.
\end{solution}

\begin{solution}{ex:k4}
\idea Bass's formula turns the count of cycles into the spectrum of a $4\times4$ matrix.
\writeup $A=J_4-I$ has eigenvalues $3$ and $-1$ (three times), and $|-1|\le2\sqrt2$: Ramanujan, with $q=2$, $n=4$, $m=6$. The eigenvalues of $B$: for $\lambda=3$, $\mu=2,1$; for $\lambda=-1$, the roots $\mu,\bar\mu$ of $\mu^2+\mu+2$, with
$\mu+\bar\mu=-1$, three times each; and $\pm1$, twice each. From $\mu^2=-\mu-2$, $\mu^3=2-\mu$ and $\mu^4=3\mu+2$, so $\mu^3+\bar\mu^3=5$ and $\mu^4+\bar\mu^4=1$. By \eqref{eq:nkspec}, $N_3=8+1+15+0=24$ and $N_4=16+1+3+4=24$. Cycles of
length $1$ and $2$ do not exist ($N_1=N_2=0$), so $3\pi(3)=N_3$ and $4\pi(4)=N_4$: $\pi(3)=8$, the $4$ triangles in $2$ directions, and $\pi(4)=6$, the $3$ Hamiltonian $4$-cycles in $2$ directions.
\end{solution}

\begin{solution}{ex:irregularbass}
\idea The regularity was used once, in $TT^{\mathsf T}=(q+1)I$.
\writeup Everything in the proof of Theorem~\ref{thm:bass} holds, except that $TT^{\mathsf T}=Q+I$ is the diagonal matrix of degrees. Then $SM^{-1}T^{\mathsf T}=\frac{A-u(Q+I)}{1-u^2}$ and
$\det(I_n-uSM^{-1}T^{\mathsf T})=(1-u^2)^{-n}\det((1-u^2)I-uA+u^2(Q+I))=(1-u^2)^{-n}\det(I-uA+u^2Q)$.
\end{solution}

\begin{solution}{ex:sl2gen}
\idea Gaussian elimination by row operations is exactly multiplication by powers of $U$ and $L$.
\writeup The first column of a matrix of $\mathrm{SL}_2(\mathbf F_p)$ is any non-zero vector ($p^2-1$ choices), and given it, the second column $v$ must satisfy the linear equation $\det(c_1,v)=1$, which has $p$ solutions. So
$|\mathrm{SL}_2(\mathbf F_p)|=p(p^2-1)$. Let $g=(\begin{smallmatrix}a&b\\c&d\end{smallmatrix})\in\mathrm{SL}_2(\mathbf F_p)$; $U^xg$ adds $x$ times the second row to the first, $L^yg$ adds $y$ times the first row to the second. If $c\ne0$, choose $x$ with
$a+xc=1$; then a multiplication by $L^y$ makes the bottom left entry $0$, and by $U^z$ the top right entry $0$, leaving a diagonal matrix of determinant $1$ whose top left entry is $1$: the identity. If $c=0$, then $a\ne0$
and one multiplication by $L$ makes $c\ne0$. So $g^{-1}$, hence $g$, is a product of powers of $U$ and $L$. The Cayley graph is connected; it is $4$-regular because $U^{\pm1}$, $L^{\pm1}$ are four distinct matrices when $p\ge3$.
\end{solution}

\section*{Sources and further reading}

\begin{itemize}
\item \emph{History.} M.~S.~Pinsker, ``On the complexity of a concentrator'', \emph{7th International Teletraffic Conference}, Stockholm, 1973, 318/1--318/4. G.~A.~Margulis, ``Explicit constructions of concentrators'', \emph{Problemy
Peredachi Informatsii} 9:4 (1973), 71--80. Y.~Ihara, ``On discrete subgroups of the two by two projective linear group over $\mathfrak p$-adic fields'', \emph{Journal of the Mathematical Society of Japan} 18 (1966), 219--235.
T.~Sunada, ``L-functions in geometry and some applications'', \emph{Lecture Notes in Mathematics} 1201 (1986), 266--284. H.~Bass, ``The Ihara--Selberg zeta function of a tree lattice'', \emph{International Journal of Mathematics} 3
(1992), 717--797. B.~D.~McKay, ``The expected eigenvalue distribution of a large regular graph'', \emph{Linear Algebra and its Applications} 40 (1981), 203--216. A.~Lubotzky, R.~Phillips and P.~Sarnak, ``Ramanujan graphs'',
\emph{Combinatorica} 8 (1988), 261--277.
\item \emph{Randomness.} M.~Ajtai, J.~Koml\'os and E.~Szemer\'edi, ``Deterministic simulation in LOGSPACE'', \emph{Proceedings of the 19th ACM Symposium on Theory of Computing} (1987), 132--140. D.~Gillman, ``A Chernoff bound for
random walks on expander graphs'', \emph{SIAM Journal on Computing} 27 (1998), 1203--1220.
\item \emph{Recent results.} A.~Marcus, D.~A.~Spielman and N.~Srivastava, ``Interlacing families I: bipartite Ramanujan graphs of all degrees'', \emph{Annals of Mathematics} 182 (2015), 307--325. J.~Friedman, \emph{A Proof of Alon's
Second Eigenvalue Conjecture and Related Problems}, Memoirs of the AMS 195, 2008. J.~Huang, T.~McKenzie and H.-T.~Yau, ``Ramanujan property and edge universality of random regular graphs'', preprint, 2024.
J.~Bourgain and A.~Gamburd, ``Uniform expansion bounds for Cayley graphs of $\mathrm{SL}_2(\mathbb F_p)$'', \emph{Annals of Mathematics} 167 (2008), 625--642. The bet: L.~Sloman, ``New proof settles decades-old bet
about connected networks'', \emph{Quanta Magazine}, 18 April 2025.
\item S.~Hoory, N.~Linial and A.~Wigderson, ``Expander graphs and their applications'', \emph{Bulletin of the AMS} 43 (2006), 439--561. The standard survey, from mixing to zig-zag products.
\item A.~Terras, \emph{Zeta Functions of Graphs: A Stroll through the Garden}, Cambridge University Press, 2011. Ihara's zeta function, its determinant formulas and graph prime number theorems.
\item G.~Davidoff, P.~Sarnak and A.~Valette, \emph{Elementary Number Theory, Group Theory and Ramanujan Graphs}, Cambridge University Press, 2003. The LPS graphs with complete proofs at the level of this series.
\end{itemize}
